% !TeX spellcheck = en_US
% !TeX encoding = UTF-8
\chapter*{Abstract}

This report documents the first lab of the Security Insider Lab. It starts with the networking background the rest of the course builds on: the TCP/IP stack, the core protocols (IP, ICMP, TCP, UDP, DNS, HTTP and TLS), the Border Gateway Protocol, the difference between Internet service providers and Internet exchange points, and DE-CIX in Frankfurt as Germany's largest exchange. Then we put these concepts into practice: using Wireshark 4.2.2 on Ubuntu 24.04 in our own home network, we captured and analyzed a ping, a DNS lookup, a TCP three-way handshake, a plain HTTP request and a TLS-protected connection. The captures matched the textbook in most respects, but a few things surprised us, such as ping running over IPv6, repeated SYN retransmissions during a connection set-up, and a client racing IPv6 against IPv4. Next, the report surveys five network attacks (DDoS, man-in-the-middle, ARP spoofing, DNS spoofing and port scanning) with a closer look at DDoS types, amplification and infrastructure defenses, reviews five attack tools from LOIC to the Aisuru/Kimwolf botnets without running them, and compares two defenses, Cloudflare and eBPF/XDP. The last part explains how machine learning can detect attacks and describes five suitable algorithms: Random Forest, XGBoost, SVM, LSTM and Isolation Forest.
